\subsection*{Assignment 13 -- Analisi delle performance degli artisti come primary vs featured}
L’Assignment 13 richiede la definizione e l’implementazione, tramite SSIS, di una query analitica di interesse business a supporto di un’attività di consulenza per una record label, basata sui risultati delle analisi precedenti e sulla fase di \textit{data understanding}.

In risposta alla consegna, è stata progettata un’analisi delle performance degli artisti distinguendo il contributo fornito come artista primario (\textit{IsPrimary = 1}) e come artista in featuring (\textit{IsPrimary = 0}). L’analisi è limitata alle canzoni pubblicate negli ultimi tre anni e, per ciascun artista e per ciascun ruolo, vengono calcolati il numero di canzoni \textit{trending} e la percentuale di canzoni \textit{trending} sul totale delle canzoni rilevanti.

I risultati sono ordinati in base alle misure ottenute, secondo l’ordine specificato dalla query analitica.

Il flusso in SSIS utilizza la \textit{FactParticipation} come sorgente principale. I dati vengono inizialmente collegati alla dimensione \textit{DimDate} tramite un \textit{Lookup}; successivamente, un \textit{Conditional Split} consente di selezionare esclusivamente le canzoni pubblicate negli ultimi tre anni, applicando la condizione:
\[
\texttt{Year >= (YEAR(GETDATE()) - 2)}
\]
Un ulteriore \textit{Lookup} sulla dimensione \textit{DimSong} permette di associare a ciascun brano la relativa categoria musicale.

Un componente \textit{Multicast} separa il ramo destinato al calcolo delle statistiche di benchmark dal ramo contenente tutte le partecipazioni artista--canzone. Il benchmark statistico viene calcolato considerando esclusivamente le canzoni con \textit{IsPrimary = 1}, evitando distorsioni dovute alla presenza di più righe per la stessa canzone generate dai featuring. Media e deviazione standard vengono calcolate tramite trasformazioni \textit{Aggregate} e \textit{Derived Column}.

I risultati del benchmark vengono successivamente combinati tramite un \textit{Merge Join}, consentendo di classificare ogni partecipazione come \textit{trending} o \textit{non trending}. I dati vengono quindi aggregati per artista e ruolo al fine di ottenere i conteggi e le percentuali richieste.

Dopo il \textit{Lookup} sulla dimensione \textit{DimArtist}, viene applicato un \textit{Conditional Split} denominato \textit{SplitOnlyArtistsVsTrending}, che seleziona esclusivamente gli artisti con almeno un brano \textit{trending}, verificando che la somma dei contributi \textit{trending} come artista primario e come artista in featuring sia maggiore di zero.

Dal punto di vista di una record label, questa analisi consente di confrontare in modo diretto le performance degli artisti nei due ruoli (\textit{primary} vs \textit{featured}), supportando decisioni strategiche legate alle collaborazioni, alla valorizzazione degli artisti e alla pianificazione delle future pubblicazioni.
